%% PART I: INTRODUCTION AND BACKGROUND
\section*{Research question and significance}

Drug discovery frequently begins with a property wish list that defines a pharmacological profile of interest, including target molecular weight, lipophilicity, polar surface area, and related descriptors. The practical challenge is to find synthesizable molecules that match it. Synthesis-aware generative models address this by assembling molecules from purchasable building blocks via validated reaction templates, guaranteeing synthesizability by construction \cite{bib1}. Several models operate in this paradigm, including SynNet \cite{bib14}, ChemProjector \cite{bib15}, SynFormer \cite{bib16}, and PrexSyn \cite{bib2}.

PrexSyn, the current state of the art among synthesis-aware generators, produces property-conditioned candidates within the Enamine REAL chemical space \cite{bib3} at high throughput ($\sim$0.26\,s per target molecule), drawing from $\sim$480,000 in-stock building blocks and 115 reaction templates. It reconstructs 94\% of molecules drawn from Enamine REAL but only 28\% from ChEMBL \cite{bib4}, a curated database of bioactive compounds with demonstrated pharmacological activity. The remaining 72\% of ChEMBL represents pharmacologically validated chemistry beyond the model's reach.

This coverage gap is structural: every synthesis-aware generator is bounded by its building block library and reaction template set. Retraining on a larger chemical space shifts the boundary but does not eliminate it, and demands substantial data curation. Repeated sampling draws from the same constrained space, producing additional candidates but opening no new chemical territory. For discovery programs seeking novel scaffolds or unexplored regions of bioactive space, this constraint directly limits the generator's practical utility.

We investigate whether post-generation molecular modification expands the space of viable candidates when composed with the frontier synthesis-aware model. Fragment substitution, matched molecular pair transforms, latent-space editing, and scaffold-constrained decoration are established modification techniques that operate outside any single building-block library and are available as open-source tools. We benchmark four such methods, two rule-based and two learned, on their ability to produce synthesizable, property-conserving candidates beyond what repeated PrexSyn sampling achieves. We stratify results by how well PrexSyn performs alone to reveal when modification adds the most marginal value. A complementarity analysis then identifies which combinations maximize total unique coverage, providing empirical guidance for pipeline design in both medicinal chemistry and automated agent-driven discovery systems.

\section*{Context and datasets}

PrexSyn \cite{bib2} generates synthesis pathways from $\sim$480,000 Enamine in-stock building blocks via 115 reaction templates, conditioned on molecular property specifications. We use it as a black-box generator that accepts a property profile and returns synthesizable candidates within Enamine REAL.

ChEMBL \cite{bib4} provides property specifications grounded in demonstrated pharmacological activity. Its molecules have been tested in real assays, making them representative of genuine drug discovery targets. From each sampled compound we compute a set of molecular property descriptors using RDKit \cite{bib18}, drawn from the physicochemical, topological, and fingerprint descriptors that PrexSyn accepts as conditioning input. These form the property specification passed to PrexSyn. Structural similarity to the ChEMBL reference then serves as a supplementary diagnostic of how close generated candidates fall to known bioactive chemistry.

AiZynthFinder \cite{bib5} serves as our synthesizability gate. It decomposes each candidate retrosynthetically into purchasable building blocks via reactions learned from the USPTO corpus \cite{bib6}. We configure it against two commercial catalogs, Enamine in-stock compounds and ZINC \cite{bib17}, extending coverage beyond PrexSyn's Enamine-only space.

The four modification methods draw on independent chemical knowledge. CReM \cite{bib7} uses a context-constrained fragment database mined from ChEMBL. mmpdb \cite{bib8} applies empirical matched molecular pair transformation rules. JT-VAE \cite{bib9} encodes and decodes molecules via a learned continuous representation. LibINVENT \cite{bib10} generates scaffold decorations via reinforcement learning, accessed through REINVENT4 \cite{bib11}.
